% Physical page 1 of the concise study, first half: the complete map of the
% appointed elements in the order of the Mass. The four overview rows follow
% immediately, in the next component, on this same page.

\zlabel{triptych:concise:inventory:start}
\begin{concisemaptable}
Introit & \latin{In voluntate tua, Domine}; Esth 13:9--11, adapted; verse
Ps 118:1 (Heb.\ 119) & Mardochai's confession under sentence of death: ``All
things, O Lord, are in thy power: and no one can resist thy will'' (1861
English). The verse blesses ``the undefiled in the way''.\\
Collect & \latin{Familiam tuam, quaesumus, Domine} & ``Preserve thy family, O
Lord, we beseech thee, by thy constant mercy'' (1861 English), free under his
protection from all adversity and devoted to his name in good works.\\
Epistle & Eph 6:10--17 & Put on the armour of God, for ``our wrestling is not
against flesh and blood; but against principalities and powers''; stand in the
evil day with the shield of faith and the sword of the Spirit.\\
Gradual & \latin{Domine, refugium}; Ps 89:1--2 (Heb.\ 90) & ``Lord, thou hast
been our refuge from generation to generation''; ``from eternity and to eternity
thou art God.''\\
Alleluia & \latin{In exitu Israel}; Ps 113:1 (Heb.\ 114) & ``When Israel went
out of Egypt, the house of Jacob from a barbarous people.''\\
Gospel & Mt 18:23--35 & A servant forgiven ten thousand talents throttles a
fellow servant for a hundred pence and is handed to the torturers: ``So also
shall my heavenly Father do to you''.\\
Offertory & \latin{Vir erat in terra Hus}; printed citation Job 1 & Job's
testing, in the chant's own words: ``Satan asked to tempt him; and power was
given him by the Lord'' (1861 English).\\
Secret & \latin{Suscipe, Domine, propitius hostias} & Offerings ``by which thou
art pleased to be appeased''; ``restore us to salvation by thy powerful
goodness'' (1861 English).\\
Communion & \latin{In salutari tuo}; Ps 118:81, 84 and 86 (Heb.\ 119), adapted
& ``When wilt thou judge those who persecute me?'' (1861 English). The question
about the number of the servant's days is not sung.\\
Postcommunion & \latin{Immortalitatis alimoniam} & ``Having received the food of
immortality \dots\ what we have taken with our mouths we may receive with pure
souls'' (1861 English).\\
\end{concisemaptable}
\zlabel{triptych:concise:inventory:end}

\vspace{0.1em}
{\footnotesize\noindent Second class, green, with Gloria, Credo and the Trinity
Preface. On 18~October 2026 St Luke is commemorated in low and conventual Masses
only, and the oration \latin{pro Fidei propagatione} is joined to the three
orations in every Mass unless the Ordinary has appointed another day for the
missions (RG 16, 108, 111, 127~b; RGMR 434, 437~a, 444~c, 445--446, 453, 475~a,
494~b).\par}
